% GENERATED FILE. Edit canonical lessons, then run npm run generate:books.
% canonical-source-hash: fnv1a64:242af8757de2f35f
% canonical-lessons: GE-C362-ehrenamt, GE-C362-ehrenamtlich, GE-C362-freiwillig, GE-C362-spenden, GE-C362-engagieren

\chapter{Freiwillig helfen}
\label{ch:ge-freiwillig-helfen}
\hlchaptermodality{362}

\begin{chapteropening}
\textbf{By the end of this chapter:} \emph{I can talk about voluntary work, donating and getting involved.}

Complete the last lesson of chapter 362: I can talk about voluntary work, donating and getting involved.
\end{chapteropening}

\section[das Ehrenamt]{das Ehrenamt — voluntary work, an unpaid post}

\label{lesson:GE-C362-ehrenamt}



\begin{quote}
Before the new one: say the German for to integrate, then the German for violence.
\end{quote}

\subsection*{You'll want to know: das Ehrenamt}
\textbf{das Ehrenamt} — \textquotedblleft{}voluntary work, an unpaid post\textquotedblright{}.

\begin{quote}
Neben dem Beruf hat sie ein Ehrenamt im Sportverein.
\end{quote}

\begin{quote}
\emph{Besides her job she has an unpaid post at the sports club.}
\end{quote}

Plural: \emph{die Ehrenämter}. Literally an office of honour.

\begin{cousinweb}
\emph{Ehrenamt} joins \emph{Ehre} \textquotedblleft{}honour\textquotedblright{}, Old High German \emph{ēra} (EtymDB), to \emph{Amt} \textquotedblleft{}office\textquotedblright{}, from Middle High German \emph{ambahte} (EtymDB), a word that goes back to Proto-Germanic \emph{ambahtaz} \textquotedblleft{}follower, servant\textquotedblright{}.
\end{cousinweb}

\begin{grammarlens}[title={What you've built}]
One more word for telling, explaining and describing at B1.
\end{grammarlens}

\subsection*{Your turn}
\begin{itemize}
\raggedright
  \item \emph{Say it:} \emph{das Ehrenamt}
  \item \emph{Say it:} \emph{das Ehrenamt}, once more
  \item \emph{Say it:} a sentence of your own with \emph{das Ehrenamt}
  \item \emph{Recall:} say the German for to integrate, then the German for violence, then say \emph{das Ehrenamt} again
\end{itemize}

\begin{tcolorbox}[breakable,colback=teal!4,colframe=teal!35!black,title={Before you move on}]
What does das Ehrenamt mean? (\textquotedblleft{}Voluntary work, an unpaid post\textquotedblright{}.) Say it once more.
\end{tcolorbox}
\section[ehrenamtlich]{ehrenamtlich — unpaid, voluntary}

\label{lesson:GE-C362-ehrenamtlich}



\begin{quote}
Before the new one: say the German for violence, then the German for voluntary work.
\end{quote}

\subsection*{You'll want to know: ehrenamtlich}
\textbf{ehrenamtlich} — \textquotedblleft{}unpaid, voluntary\textquotedblright{}.

\begin{quote}
Mein Großvater arbeitet ehrenamtlich in der Bibliothek.
\end{quote}

\begin{quote}
\emph{My grandfather works as a volunteer in the library.}
\end{quote}

\emph{Ehrenamtlich arbeiten} — to do voluntary work.

\begin{cousinweb}
\emph{ehrenamtlich} is built with \emph{-lich} on \emph{Ehrenamt}, from \emph{Ehre} \textquotedblleft{}honour\textquotedblright{}, Old High German \emph{ēra}, and \emph{Amt} \textquotedblleft{}office\textquotedblright{}, Middle High German \emph{ambahte} (EtymDB).
\end{cousinweb}

\begin{grammarlens}[title={What you've built}]
One more word for telling, explaining and describing at B1.
\end{grammarlens}

\subsection*{Your turn}
\begin{itemize}
\raggedright
  \item \emph{Say it:} \emph{ehrenamtlich}
  \item \emph{Say it:} \emph{ehrenamtlich}, once more
  \item \emph{Say it:} a sentence of your own with \emph{ehrenamtlich}
  \item \emph{Recall:} say the German for violence, then the German for voluntary work, then say \emph{ehrenamtlich} again
\end{itemize}

\begin{tcolorbox}[breakable,colback=teal!4,colframe=teal!35!black,title={Before you move on}]
What does ehrenamtlich mean? (\textquotedblleft{}Unpaid, voluntary\textquotedblright{}.) Say it once more.
\end{tcolorbox}
\section[freiwillig]{freiwillig — voluntary, of your own free will}

\label{lesson:GE-C362-freiwillig}



\begin{quote}
Before the new one: say the German for voluntary work, then the German for unpaid.
\end{quote}

\subsection*{You'll want to know: freiwillig}
\textbf{freiwillig} — \textquotedblleft{}voluntary, of your own free will\textquotedblright{}.

\begin{quote}
Niemand hat ihn gezwungen, er hat es freiwillig gemacht.
\end{quote}

\begin{quote}
\emph{Nobody forced him; he did it of his own free will.}
\end{quote}

\emph{Der Freiwillige, die Freiwillige} — the volunteer.

\begin{cousinweb}
\emph{freiwillig} joins \emph{frei} \textquotedblleft{}free\textquotedblright{}, Proto-Germanic \emph{frijaz} (EtymDB), to \emph{Wille} \textquotedblleft{}will\textquotedblright{}, with \emph{-ig}: with a free will.
\end{cousinweb}

\begin{grammarlens}[title={What you've built}]
One more word for telling, explaining and describing at B1.
\end{grammarlens}

\subsection*{Your turn}
\begin{itemize}
\raggedright
  \item \emph{Say it:} \emph{freiwillig}
  \item \emph{Say it:} \emph{freiwillig}, once more
  \item \emph{Say it:} a sentence of your own with \emph{freiwillig}
  \item \emph{Recall:} say the German for voluntary work, then the German for unpaid, then say \emph{freiwillig} again
\end{itemize}

\begin{tcolorbox}[breakable,colback=teal!4,colframe=teal!35!black,title={Before you move on}]
What does freiwillig mean? (\textquotedblleft{}Voluntary, of your own free will\textquotedblright{}.) Say it once more.
\end{tcolorbox}
\section[spenden]{spenden — to donate, to give}

\label{lesson:GE-C362-spenden}



\begin{quote}
Before the new one: say the German for unpaid, then the German for voluntary.
\end{quote}

\subsection*{You'll want to know: spenden}
\textbf{spenden} — \textquotedblleft{}to donate, to give\textquotedblright{}.

\begin{quote}
Jedes Jahr spenden wir Geld für Kinder in Not.
\end{quote}

\begin{quote}
\emph{Every year we donate money for children in need.}
\end{quote}

\emph{Blut spenden} — to give blood. \emph{Die Spende} — the donation.

\begin{cousinweb}
\emph{spenden} goes back to Middle High German \emph{spenden} and Old High German \emph{spentōn} \textquotedblleft{}to spend, to use\textquotedblright{} (EtymDB). It is usually traced to Medieval Latin \emph{spendere}, from \emph{expendere} \textquotedblleft{}to weigh out, to pay out\textquotedblright{} (Lewis \& Short); English \emph{spend} is a cousin.
\end{cousinweb}

\begin{grammarlens}[title={What you've built}]
One more word for telling, explaining and describing at B1.
\end{grammarlens}

\subsection*{Your turn}
\begin{itemize}
\raggedright
  \item \emph{Say it:} \emph{spenden}
  \item \emph{Say it:} \emph{spenden}, once more
  \item \emph{Say it:} a sentence of your own with \emph{spenden}
  \item \emph{Recall:} say the German for unpaid, then the German for voluntary, then say \emph{spenden} again
\end{itemize}

\begin{tcolorbox}[breakable,colback=teal!4,colframe=teal!35!black,title={Before you move on}]
What does spenden mean? (\textquotedblleft{}To donate, to give\textquotedblright{}.) Say it once more.
\end{tcolorbox}
\section[engagieren]{engagieren — to get involved (mostly reflexive, sich engagieren); to hire (an artist)}

\label{lesson:GE-C362-engagieren}



\begin{quote}
Before the new one: say the German for voluntary, then the German for to donate.
\end{quote}

\subsection*{You'll want to know: engagieren}
\textbf{engagieren} — \textquotedblleft{}to get involved (mostly reflexive, sich engagieren); to hire (an artist)\textquotedblright{}.

\begin{quote}
Sie engagiert sich seit Jahren für den Tierschutz.
\end{quote}

\begin{quote}
\emph{She has been involved in animal protection for years.}
\end{quote}

Mostly reflexive: \emph{sich für etwas engagieren}. Said the French way: \emph{an-ga-ZHEE-ren}. \emph{Einen Musiker engagieren} — to hire a musician.

\begin{cousinweb}
\emph{engagieren} was borrowed from French \emph{engager} \textquotedblleft{}to pledge, to bind\textquotedblright{} (EtymDB), built on \emph{gage} \textquotedblleft{}a pledge\textquotedblright{}, which is usually traced to a Frankish word; EtymDB lists Old High German \emph{wetti} \textquotedblleft{}a pledge, a security\textquotedblright{} as a cousin.
\end{cousinweb}

\begin{grammarlens}[title={What you've built}]
One more word for telling, explaining and describing at B1.
\end{grammarlens}

\subsection*{Your turn}
\begin{itemize}
\raggedright
  \item \emph{Say it:} \emph{engagieren}
  \item \emph{Say it:} \emph{engagieren}, once more
  \item \emph{Say it:} a sentence of your own with \emph{engagieren}
  \item \emph{Recall:} say the German for voluntary, then the German for to donate, then say \emph{engagieren} again
\end{itemize}

\begin{tcolorbox}[breakable,colback=teal!4,colframe=teal!35!black,title={Before you move on}]
What does engagieren mean? (\textquotedblleft{}To get involved (mostly reflexive, sich engagieren); to hire (an artist)\textquotedblright{}.) Say it once more.
\end{tcolorbox}
